\chapter{EMA 锚点与 OPD 的联合动力学}
\label{chap:ema}

\section{两个策略状态、三种时钟}

hard sync 在 outer 边界把训练后的 actor 完整复制为下一轮 reference。指数移动平均（exponential moving average，EMA）保存一个独立的参数状态，因此 actor、rollout policy 和 KL anchor 必须分别记号。设 actor 参数为 $\theta_t$，EMA 参数为 $\phi_t$，固定 teacher 为 $T$，并写
\[
 p_t=\pi_{\theta_t},\qquad q_t=\pi_{\phi_t},\qquad
 \delta_t=\theta_t-\phi_t.
\]
rollout policy 是 batch 开始时保存的 actor 副本 $\pi_{\mathrm{old}}$。它用于采样和重要性比率；$q_t$ 用于 KL 正则化。二者具有不同的更新时钟。

本章保留 actor 学习率 $\alpha>0$，额外 EMA--KL 系数 $\beta\ge0$，EMA 保留系数 $h=\eta_{\mathrm{ema}}\in[0,1)$，MirrorOPD 步长 $\eta=\eta_{\mathrm{opd}}\in(0,1]$。$h$ 是旧 anchor 的保留比例，$1-h$ 是新 actor 的混入比例。$T_{\mathrm{ema}}$ 表示两次 anchor 更新之间的 rollout iteration 数；每个 iteration 的 inner update 数记为 $\Lambda$。

\subsection{原论文的理论时序与实现时序}

EMA Policy Gradient 论文的局部动力学使用更新前 actor：\cite{ema2026}
\begin{align}
 \theta_{t+1}&=\theta_t+\alpha\left[g_t-
 \beta\nabla_\theta\KL(\pi_\theta\Vert\pi_{\phi_t})\big|_{\theta_t}\right],
 \label{eq:ema-pre-actor}\\
 \phi_{t+1}&=h\phi_t+(1-h)\theta_t.
 \label{eq:ema-pre-anchor}
\end{align}
这里 $t$ 是 actor step 的时钟。推导滞后递推时必须保留 $\theta_t$，否则稳定矩阵会改变。

论文算法框先采样 rollout，再执行 $\Lambda$ 次 actor 更新，在 $t\bmod T_{\mathrm{ema}}=0$ 时用已经更新后的 actor 做一次 EMA。这里 $t$ 是 rollout iteration 的时钟。若每轮 $\Lambda$ 固定，两个 EMA 事件之间共有 $N=T_{\mathrm{ema}}\Lambda$ 次 actor 更新。算法框与单步理论方程的差别包括同步时机和同步频率，后文分别推导。

\subsection{替换 reference 与新增 KL}

将 MirrorOPD reference 替换为 $q$，定义固定 anchor 的奖励与基础目标
\begin{align}
 r_q(a)&=\log T(a)-\log q(a),\\
 \Psi_q(p)&=\E_{a\sim p}[r_q(a)]-\eta^{-1}\KL(p\Vert q).
 \label{eq:ema-base-objective}
\end{align}
只做这种替换时，最优策略是 $G_\eta(q)$。若在基础目标之外新增 $-\beta\KL(p\Vert q)$，总目标变为
\begin{equation}
 \Psi_{q,\beta}(p)=\E_p[r_q]-c\KL(p\Vert q),
 \qquad c=\eta^{-1}+\beta.
 \label{eq:ema-total-objective}
\end{equation}
新增 KL 改变总目标的最优解，不能解释为仅减缓有限 inner steps，却保持 exact target 不变。

\begin{theorem}[额外同 anchor KL 的 exact target]
\label{thm:ema-effective-eta}
在有限动作集上假设 $q(a)>0$、$T(a)>0$。令
\[
 \gamma=\frac1c=\frac{\eta}{1+\beta\eta}.
\]
目标\eqref{eq:ema-total-objective}的唯一最优解与最优性缺口为
\begin{align}
 p^*_{q,\beta}(a)&=\frac{q(a)^{1-\gamma}T(a)^\gamma}{Z_\gamma(q)}
 =G_\gamma(q)(a),\\
 Z_\gamma(q)&=\sum_bq(b)^{1-\gamma}T(b)^\gamma,\\
 \Psi_{q,\beta}(p^*_{q,\beta})-\Psi_{q,\beta}(p)
 &=c\KL(p\Vert p^*_{q,\beta}).
\end{align}
\end{theorem}
\begin{proof}
将 $\log p^*=(1-\gamma)\log q+\gamma\log T-\log Z_\gamma$ 代入 KL：
\begin{align*}
 -c\KL(p\Vert p^*)
 &=-c\sum_ap(a)\log p(a)
   +c(1-\gamma)\sum_ap(a)\log q(a)\\
 &\quad+c\gamma\sum_ap(a)\log T(a)-c\log Z_\gamma\\
 &=-c\sum_ap(a)\log p(a)
   +(c-1)\sum_ap(a)\log q(a)\\
 &\quad+\sum_ap(a)\log T(a)-c\log Z_\gamma\\
 &=\Psi_{q,\beta}(p)-c\log Z_\gamma.
\end{align*}
所以 $\Psi_{q,\beta}(p)=c\log Z_\gamma-c\KL(p\Vert p^*)$。KL 非负且仅在 $p=p^*$ 时为零，给出唯一最优解及缺口公式。
\end{proof}

当 $\eta=1$、$\beta=0$ 时，fixed-anchor exact target 是 $T$。当 $\eta=1$、$\beta>0$ 时，$\gamma=1/(1+\beta)<1$，target 是 $q$ 与 $T$ 的几何插值。

如果 MirrorOPD reference 为 $q$，新增 EMA--KL 指向另一个 anchor $u$，总目标为
\[
 \E_p\log(T/q)-\eta^{-1}\KL(p\Vert q)-\beta\KL(p\Vert u).
\]
按同样的展开，其 exact target 为
\[
 p^*(a)\ \propto\ T(a)^{1/c}
 q(a)^{(\eta^{-1}-1)/c}u(a)^{\beta/c}.
\]
两个 anchor 的区分直接进入 target。若新增项使用 forward KL $\KL(u\Vert p)$，导数包含 $-u(a)/p(a)$，上述几何闭式解不再适用。

\section{准稳态滞后定理及其边界}

\subsection{Fisher 近似}

设策略族在当前邻域三阶可微，固定状态测度下的 Fisher 矩阵为
\[
 F_\theta=\E_s\E_{a\sim\pi_\theta(\cdot\mid s)}
 [\nabla_\theta\log\pi_\theta(a\mid s)
  \nabla_\theta\log\pi_\theta(a\mid s)^\top].
\]
固定 $\phi$ 求导的局部展开是
\begin{align}
 \KL(\pi_\theta\Vert\pi_\phi)
 &=\tfrac12(\theta-\phi)^\top F(\theta-\phi)
   +O(\|\theta-\phi\|^3),\\
 \nabla_\theta\KL(\pi_\theta\Vert\pi_\phi)
 &=F(\theta-\phi)+O(\|\theta-\phi\|^2).
\end{align}
下面的线性递推是冻结 Fisher 后局部模型的精确结论。非线性策略只有在轨迹留在展开邻域、余项受控时才继承局部结论。

\begin{theorem}[更新前 EMA 的滞后递推与准稳态闭式]
\label{thm:ema-lag}
在线性模型\eqref{eq:ema-pre-actor}--\eqref{eq:ema-pre-anchor}中，
\[
 \delta_{t+1}=(hI-\alpha\beta F_t)\delta_t+\alpha g_t.
\]
如果窗口内 $F_t=F$、$g_t=g$，定义
\[
 D=hI-\alpha\beta F,\qquad
 S_k=\sum_{j=0}^{k-1}D^j,\qquad U_k=\sum_{j=0}^{k-1}S_j.
\]
则
\begin{align}
 \delta_{t+k}&=D^k\delta_t+\alpha S_kg,\\
 \theta_{t+k}&=\theta_t-\alpha\beta FS_k\delta_t
       +(\alpha kI-\alpha^2\beta FU_k)g.
 \label{eq:ema-quasi-closed}
\end{align}
由于 $I-D=(1-h)I+\alpha\beta F\succ0$，还成立
\[
 S_k=(I-D^k)(I-D)^{-1},\qquad
 U_k=(kI-S_k)(I-D)^{-1}.
\]
\end{theorem}
\begin{proof}
用 actor 更新减去 anchor 更新：
\begin{align*}
 \delta_{t+1}
 &=\theta_t+\alpha g_t-\alpha\beta F_t\delta_t
   -h\phi_t-(1-h)\theta_t\\
 &=h\delta_t-\alpha\beta F_t\delta_t+\alpha g_t.
\end{align*}
冻结系数后展开 $k$ 次，得到 $D^k\delta_t+\alpha\sum_{j=0}^{k-1}D^jg$。再将 actor 增量求和：
\begin{align*}
 \theta_{t+k}-\theta_t
 &=\sum_{j=0}^{k-1}(\alpha g-\alpha\beta F\delta_{t+j})\\
 &=\alpha kg-\alpha\beta F\sum_{j=0}^{k-1}
   (D^j\delta_t+\alpha S_jg),
\end{align*}
给出\eqref{eq:ema-quasi-closed}。等比求和恒等式 $(I-D)S_k=I-D^k$ 给出 $S_k$，再次求和给出
\[
 U_k=\sum_{j=0}^{k-1}(I-D^j)(I-D)^{-1}
     =(kI-S_k)(I-D)^{-1}.
\]
$F\succeq0$、$h<1$ 保证 $I-D\succ0$，所以逆矩阵存在。
\end{proof}

\begin{theorem}[滞后稳定性、振荡与稳态速度]
\label{thm:ema-lag-stability}
设 $F\succeq0$，$\lambda_i$ 为其特征值，$\chi_i=h-\alpha\beta\lambda_i$。则
\[
 \delta_{t+k}^{(i)}=\chi_i^k\delta_t^{(i)}
 +\alpha\frac{1-\chi_i^k}{1-\chi_i}g^{(i)}.
\]
滞后相对于准稳态渐近稳定的充要条件为
\[
 \alpha\beta\lambda_{\max}<1+h.
\]
稳定时，滞后极限与 actor 渐近增量是
\begin{align}
 \delta_*^{(i)}&=\frac{\alpha g^{(i)}}{1-h+\alpha\beta\lambda_i},\\
 \lim_{k\to\infty}(\theta_{t+k+1}^{(i)}-\theta_{t+k}^{(i)})
 &=\frac{\alpha(1-h)}{1-h+\alpha\beta\lambda_i}g^{(i)}.
\end{align}
\end{theorem}
\begin{proof}
投影到 $F$ 的正交特征基，得到标量递推 $d_{j+1}=\chi_i d_j+\alpha g^{(i)}$，等比求和给出闭式。瞬态衰减等价于 $|\chi_i|<1$。因 $\chi_i\le h<1$，只需 $h-\alpha\beta\lambda_i>-1$；对所有方向取最大特征值得到条件。

当 $0\le\alpha\beta\lambda_i\le h$ 时，$\chi_i\in[0,1)$，与稳态的偏差不发生符号翻转。当 $h<\alpha\beta\lambda_i<1+h$ 时，$\chi_i\in(-1,0)$，偏差交替衰减。边界上乘子为 $-1$，非零瞬态保持幅度；超过边界时幅度增长。边界应称为失去渐近稳定，而非一律指数发散。

稳定时令 $k\to\infty$ 得到 $\delta_*^{(i)}$，代入 actor 的 $\alpha g^{(i)}-\alpha\beta\lambda_i\delta_*^{(i)}$，合并分母得到速度公式。
\end{proof}

该定理允许 $\beta=0$ 和 Fisher 零特征值，此时 $\chi_i=h$，滞后仍然稳定。稳定条件没有额外的严格正下界。固定非零 $g$ 时 actor 仍以非零速度前进，所以 lag 稳定不等价于 actor 收敛，更不等价于到达 teacher。

准稳态公式给出 EMA 减慢更新的准确条件：固定 $g$ 时，$\beta\lambda_i>0$ 的 actor 方向，其渐近速度乘以
\[
 \frac{1-h}{1-h+\alpha\beta\lambda_i}<1.
\]
这是固定梯度局部模型的速度比较；梯度噪声、允许步长和曲率随策略改变时，不能据此宣称 EMA 在所有训练预算下都更慢。

\section{teacher 误差与滞后的联合收敛}

\subsection{固定局部目标}

设可表示目标参数为 $\theta^*$，actor 误差 $e_t=\theta_t-\theta^*$。若基础上升梯度在目标附近为 $g_t=-Ae_t+\xi_t$，$A\succ0$，更新前 EMA 联合系统是
\begin{equation}
 z_{t+1}=Mz_t+\alpha B\xi_t,
 \quad z_t=\begin{bmatrix}e_t\\\delta_t\end{bmatrix},\quad
 B=\begin{bmatrix}I\\I\end{bmatrix},\quad
 M=\begin{bmatrix}I-\alpha A&-\alpha\beta F\\
                  -\alpha A&hI-\alpha\beta F\end{bmatrix}.
 \label{eq:ema-joint-fixed}
\end{equation}
上升梯度的单调性符号为 $\langle g(e),e\rangle\le-m\|e\|^2$；损失下降梯度才使用正号。

\begin{theorem}[共同特征模态的精确 Jury 条件]
\label{thm:ema-jury}
假设 $A,F$ 对称可交换，$A\succ0$、$F\succeq0$。在共同特征方向写
\[
 a_i=\alpha\lambda_i(A)>0,\qquad b_i=\alpha\beta\lambda_i(F)\ge0.
\]
系统\eqref{eq:ema-joint-fixed}渐近稳定的充要条件是每个方向满足
\begin{equation}
 a_i+\frac{2b_i}{1+h}<2.
 \label{eq:ema-jury-simple}
\end{equation}
\end{theorem}
\begin{proof}
对称可交换矩阵可同时正交对角化，第 $i$ 个方向的矩阵为
\[
 M_i=\begin{bmatrix}1-a_i&-b_i\\-a_i&h-b_i\end{bmatrix}.
\]
迹与行列式分别为
\[
 \tau_i=1+h-a_i-b_i,\qquad
 d_i=(1-a_i)(h-b_i)-a_ib_i=h(1-a_i)-b_i.
\]
特征多项式是 $x^2-\tau_ix+d_i$。实二阶多项式根严格位于单位圆内的 Jury 条件是
\[
 1-\tau_i+d_i>0,\quad1+\tau_i+d_i>0,\quad1-d_i>0.
\]
可直接验证这一条件：实根时前两式分别排除穿越 $+1$、$-1$，并结合 $d_i<1$ 排除两根同时位于圆外；共轭根时 $d_i$ 是模的平方，第三式控制其模。逐项代入得
\begin{align*}
 1-\tau_i+d_i&=a_i(1-h)>0,\\
 1+\tau_i+d_i&=2(1+h)-(1+h)a_i-2b_i,\\
 1-d_i&=1-h+ha_i+b_i>0.
\end{align*}
第一、第三式自动成立；第二式除以 $1+h$ 恰为\eqref{eq:ema-jury-simple}。
\end{proof}

当 $A,F$ 不可交换时，应检查完整块矩阵的谱半径，逐模态条件不再是充要条件。仅检查 lag 会漏掉 teacher 误差到 lag 的反馈。例如 $\beta=0$ 时 lag 条件总成立，actor 仍要求 $0<\alpha\lambda_i(A)<2$。

\subsection{移动 MirrorOPD anchor}

固定目标模型与 MirrorOPD 的移动 anchor 不同：$G_\gamma(q_t)$ 随 $q_t$ 移动。在 teacher 附近取可辨识局部坐标，令 $v_t=\phi_t-\theta^*$、$e_t=\theta_t-\theta^*$，冻结共同 Fisher 为 $F$。利用
\[
 \E_p\log(T/q)=\KL(p\Vert q)-\KL(p\Vert T),
\]
实际总目标精确改写为
\[
 \Psi_{q_t,\beta}(p)=-\KL(p\Vert T)-(c-1)\KL(p\Vert q_t).
\]
局部二次展开给出
\[
 \Psi_{q_t,\beta}(\pi_{\theta_t})
 \simeq-\tfrac12e_t^\top Fe_t-
 \tfrac{c-1}{2}(e_t-v_t)^\top F(e_t-v_t).
\]
因此总梯度为 $-F[e_t+(c-1)\delta_t]$，联合系统为
\begin{equation}
 \begin{bmatrix}e_{t+1}\\\delta_{t+1}\end{bmatrix}
 =\begin{bmatrix}I-\alpha F&-\alpha(c-1)F\\
                  -\alpha F&hI-\alpha(c-1)F\end{bmatrix}
   \begin{bmatrix}e_t\\\delta_t\end{bmatrix}.
 \label{eq:ema-joint-opd}
\end{equation}

\begin{corollary}[EMA--MirrorOPD 局部 teacher 收敛条件]
\label{cor:ema-opd-jury}
在上述二次模型中，若 $F\succ0$、$\eta\in(0,1]$、$\beta\ge0$，则\eqref{eq:ema-joint-opd}收敛到 $(e,\delta)=(0,0)$ 的充要条件为
\[
 0<\alpha\lambda_i(F)<\frac{2(1+h)}{2c+h-1}\quad\text{对所有 }i,
 \qquad c=\eta^{-1}+\beta.
\]
\end{corollary}
\begin{proof}
$c-1\ge0$。在定理\ref{thm:ema-jury}中取 $a_i=\alpha\lambda_i(F)$、$b_i=\alpha(c-1)\lambda_i(F)$，得到
\[
 \alpha\lambda_i(F)\left(1+\frac{2(c-1)}{1+h}\right)<2.
\]
括号合并为 $(2c+h-1)/(1+h)$，整理即得。若 $F$ 奇异，结论只作用于正特征值对应的策略变化方向，零特征方向的参数可保持初值。
\end{proof}

使用冻结 $F^{-1}$ 的 Natural GD 时，所有正 Fisher 模态的 $\alpha\lambda_i(F)$ 替换为 $\alpha$，消除局部条件数对步长上界的影响，但保留 $c,h$ 与同步时序的耦合。

\subsection{非正规矩阵与扰动速率}

\begin{theorem}[稳定联合系统的扰动界]
\label{thm:ema-perturbation}
设固定矩阵 $M$ 满足 $\rho(M)<1$，$z_{t+1}=Mz_t+b_t$。对任意 $r\in(\rho(M),1)$，存在 $C_r<\infty$，使
\[
 \|z_k\|\le C_rr^k\|z_0\|
 +C_r\sum_{j=0}^{k-1}r^{k-1-j}\|b_j\|.
\]
若 $\limsup_j\|b_j\|\le\bar b$，则 $\limsup_k\|z_k\|\le C_r\bar b/(1-r)$。
\end{theorem}
\begin{proof}
递推展开为 $z_k=M^kz_0+\sum_{j=0}^{k-1}M^{k-1-j}b_j$。Jordan 分解中，大小为 $\ell$ 的块产生至多 $k^{\ell-1}|\lambda|^{k-\ell+1}$ 的项。因 $|\lambda|<r$，几何衰减吸收有限次多项式，存在 $C_r$ 使 $\|M^k\|\le C_rr^k$，三角不等式给出首式。

给定 $\varepsilon>0$，取 $J$ 使 $j\ge J$ 时 $\|b_j\|\le\bar b+\varepsilon$。有限项 $j<J$ 随 $k$ 衰减为零，其余求和至多 $(\bar b+\varepsilon)/(1-r)$。取上极限，再令 $\varepsilon\downarrow0$ 得到第二式。
\end{proof}

一般不能把 $r$ 改成 $\rho(M)$。例如 $M=\begin{bmatrix}\rho&1\\0&\rho\end{bmatrix}$，$0<\rho<1$ 时，$M^k$ 的右上角为 $k\rho^{k-1}$，不存在固定 $C$ 使所有 $k$ 满足 $\|M^k\|\le C\rho^k$。可对角化或正规矩阵具有更强的界，但需要列出假设。

若局部目标 $\vartheta_t$ 漂移，令 $e_t=\theta_t-\vartheta_t$，额外输入为
\[
 b_t=\alpha\begin{bmatrix}I\\I\end{bmatrix}\xi_t
 +\begin{bmatrix}\vartheta_t-\vartheta_{t+1}\\0\end{bmatrix}.
\]
目标漂移只直接进入 $e$ 方程，$\delta$ 中的目标平移相消。定理\ref{thm:ema-perturbation}分别控制采样误差、有限 inner 求解误差和目标漂移。

\section{延迟 EMA 的事件级复合算子}

\subsection{固定 anchor 窗口与同步时序}

一个 anchor 事件之间 $\phi$ 固定。先冻结 $F,g$，令 $L=I-\alpha\beta F$，窗口长度为 $N$ 次 actor 更新，则
\[
 \delta_j^-=L^j\delta_0+\alpha\sum_{i=0}^{j-1}L^ig,
 \qquad\theta_j=\phi+\delta_j^-,\quad0\le j\le N.
\]
$h$ 在窗口内部没有出现。若末端使用更新后 actor，则
\begin{align}
 \phi^+&=h\phi+(1-h)\theta_N,\\
 \delta^+&=h\delta_N^-
 =hL^N\delta_0+h\alpha\sum_{i=0}^{N-1}L^ig.
 \label{eq:ema-block-post}
\end{align}
若使用最后一次 actor 更新前的参数，则
\begin{align}
 \phi^+&=h\phi+(1-h)\theta_{N-1},\\
 \delta^+&=\delta_N^--(1-h)\delta_{N-1}^-
 =(hI-\alpha\beta F)\delta_{N-1}^-+\alpha g.
 \label{eq:ema-block-pre}
\end{align}
这些公式直接由 $\delta^+=\theta_N-\phi^+$ 得到。$N=1$ 时，两种时序分别产生 $h(I-\alpha\beta F)$ 与 $hI-\alpha\beta F$，稳定乘子已经不同。

在论文算法中，$N=T_{\mathrm{ema}}\Lambda$ 仅在每轮 inner 数固定时成立。rollout 刷新改变 $g$ 和状态测度，固定系数闭式应替换为逐步矩阵乘积与输入求和。延迟更新只在事件边界应用一次 $h$，不能替换为逐 actor step EMA 中的 $h^N$。

\subsection{完整 teacher 误差的 block 矩阵}

对同一局部 MirrorOPD 总目标，窗口内状态 $(e,v)$ 满足
\[
 \begin{bmatrix}e_{j+1}\\v_{j+1}\end{bmatrix}
 =K\begin{bmatrix}e_j\\v_j\end{bmatrix},\qquad
 K=\begin{bmatrix}I-\alpha cF&\alpha(c-1)F\\0&I\end{bmatrix}.
\]
更新后同步矩阵为
\[
 S=\begin{bmatrix}I&0\\(1-h)I&hI\end{bmatrix},
 \qquad M_{\mathrm{block}}=SK^N.
\]
因此包含全部参数的事件级判据是 $\rho(SK^{T_{\mathrm{ema}}\Lambda})<1$。在共同 Fisher 模态上记
\[
 \ell_i=1-\alpha c\lambda_i,\qquad
 w_i=\frac{c-1}{c}(1-\ell_i^N).
\]
等比求和给出
\[
 K_i^N=\begin{bmatrix}\ell_i^N&w_i\\0&1\end{bmatrix},\qquad
 M_{\mathrm{block},i}=\begin{bmatrix}\ell_i^N&w_i\\
 (1-h)\ell_i^N&h+(1-h)w_i\end{bmatrix}.
\]
将迹 $\ell_i^N+h+(1-h)w_i$ 与行列式 $h\ell_i^N$ 代入 Jury 三条件即可检查。该条件同时包含 actor 窗口和同步事件；单独的 lag 条件无法替代。

\section{具有全局定理的几何 soft sync}

\subsection{更新后几何同步}

在概率的对数坐标定义几何同步
\[
 \mathcal S_h(q,p)(a)=\frac{q(a)^hp(a)^{1-h}}
 {\sum_bq(b)^hp(b)^{1-h}}.
\]
它对 logits 或 log-probability 作线性插值再归一化。定义
\[
 d_{\mathrm{proj}}(p,T)=\osc_a\log\frac{p(a)}{T(a)},
 \qquad\osc f=\max_af(a)-\min_af(a).
\]

\begin{theorem}[exact actor 后几何同步的外层收缩]
\label{thm:ema-geometric-exact}
在正有限动作分布上执行
\[
 p_{k+1}=G_\gamma(q_k),\qquad
 q_{k+1}=\mathcal S_h(q_k,p_{k+1}),\qquad
 \gamma=\frac{\eta}{1+\beta\eta}.
\]
则
\begin{align}
 q_{k+1}&=G_{(1-h)\gamma}(q_k),\\
 d_{\mathrm{proj}}(q_{k+1},T)
 &=[1-(1-h)\gamma]d_{\mathrm{proj}}(q_k,T),\\
 d_{\mathrm{proj}}(p_{k+1},T)
 &=(1-\gamma)d_{\mathrm{proj}}(q_k,T).
\end{align}
\end{theorem}
\begin{proof}
代入 $p_{k+1}\propto q_k^{1-\gamma}T^\gamma$，归一化常数相消，得到
\[
 q_{k+1}\propto q_k^h(q_k^{1-\gamma}T^\gamma)^{1-h}
 =q_k^{1-(1-h)\gamma}T^{(1-h)\gamma}.
\]
对任意 $\tau\in[0,1]$，
\[
 \log\frac{G_\tau(q)(a)}{T(a)}
 =(1-\tau)\log\frac{q(a)}{T(a)}-\log Z_\tau(q).
\]
振幅消去归一化常数，得到 $d_{\mathrm{proj}}(G_\tau(q),T)=(1-\tau)d_{\mathrm{proj}}(q,T)$。分别取 $\tau=(1-h)\gamma$、$\tau=\gamma$ 即得。
\end{proof}

hard sync 是 $h=0$，anchor 收缩因子为 $1-\gamma$。固定 $\gamma>0$ 时，$h>0$ 严格增大该因子，因此在 exact、同等 outer 预算下 anchor 的确学得更慢。$\gamma=1$ 时 actor 每轮直接等于 teacher，anchor 仍以因子 $h$ 接近 teacher；评价 actor 与评价 anchor 会得到不同答案。

有限 inner solve 的对数误差写成
\[
 \log p_{k+1}(a)=\log G_\gamma(q_k)(a)
 +\varepsilon_{k+1}(a)-C_{k+1}.
\]
若 $\osc\varepsilon_{k+1}\le e_{k+1}$，代入几何同步并取振幅得到
\[
 d_{\mathrm{proj}}(q_{k+1},T)
 \le[1-(1-h)\gamma]d_{\mathrm{proj}}(q_k,T)+(1-h)e_{k+1}.
\]
当 $e_k\le\bar e$ 时，等比求和给出距离上极限 $\bar e/\gamma$。soft sync 同时减弱理想推进与求解误差注入，故在此确定性最坏情形界里 $h$ 不改善误差地板。

\subsection{更新前几何同步的一步延迟}

若 $p_{k+1}=G_\gamma(q_k)$，但同步用 $p_k$，即 $q_{k+1}=\mathcal S_h(q_k,p_k)$，选基准动作 $a_0$ 定义
\[
 x_k(a)=\log\frac{p_k(a)T(a_0)}{p_k(a_0)T(a)},\qquad
 y_k(a)=\log\frac{q_k(a)T(a_0)}{q_k(a_0)T(a)}.
\]
归一化常数消去，每个动作方向满足
\[
 \begin{bmatrix}x_{k+1}\\y_{k+1}\end{bmatrix}
 =\begin{bmatrix}0&1-\gamma\\1-h&h\end{bmatrix}
 \begin{bmatrix}x_k\\y_k\end{bmatrix}.
\]
其特征根为
\[
 r_\pm=\frac{h\pm\sqrt{h^2+4(1-h)(1-\gamma)}}2.
\]
$0<\gamma\le1$、$h<1$ 时，$r_+<1$ 且 $|r_-|\le r_+$：判别式平方根小于 $2-h$ 恰好等价于 $\gamma>0$。故系统收敛，负根解释一步延迟的交替分量。$h=0,\gamma=1$ 时矩阵平方为零，有限两步归零。此速率区别于更新后同步的 $1-(1-h)\gamma$。

\subsection{参数 EMA 的非等价性}

参数 EMA 是 $\phi^+=h\phi+(1-h)\theta$。只有固定状态 logits 对参数仿射时，才有
\[
 \pi_{\phi^+}=\mathcal S_h(\pi_\phi,\pi_\theta).
\]
神经网络 logits 一般不对所有权重仿射。反例取二分类策略 $\pi_\theta(1)=\operatorname{sigmoid}(\theta^2)$，$\phi=-1$、$\theta=1$、$h=1/2$。端点策略相同，几何同步保持它；参数 EMA 给出 $\phi^+=0$，概率为 $1/2$，不同于端点的 $\operatorname{sigmoid}(1)$。

若 logits 映射二阶可微且 $\theta-\phi$ 小，Taylor 展开给出
\[
 z_{h\phi+(1-h)\theta}=hz_\phi+(1-h)z_\theta
 +O(\|\theta-\phi\|^2).
\]
二者在局部一阶对应，精确 projective 定理属于几何同步。二阶偏差可作为扰动控制，不能直接删去。

\section{Top-$k$ KL 与噪声动力学}

\subsection{无偏尾部校正}

EMA-PG 在指定 token 集合 $A$ 上精确计算部分 KL，再对尾部采样校正。固定状态、参数与集合，记 $p=\pi_\theta$、$q=\pi_\phi$、采样分布 $\mu$。假设 $p(a)>0$ 时 $\mu(a)>0$，reverse KL 值估计为
\[
 \widehat K=\sum_{a\in A}p(a)\log\frac{p(a)}{q(a)}
 +\1\{Y\notin A\}\frac{p(Y)}{\mu(Y)}\log\frac{p(Y)}{q(Y)},
 \quad Y\sim\mu.
\]
逐项取期望，得到
\[
 \E_\mu\widehat K=\sum_{a\in A}p(a)\log\frac{p(a)}{q(a)}
 +\sum_{a\notin A}p(a)\log\frac{p(a)}{q(a)}=\KL(p\Vert q).
\]
梯度还需要正确停止梯度与 score-function 项。真实梯度为
\begin{align*}
 \nabla_\theta\KL(p\Vert q)
 &=\sum_ap(a)\nabla_\theta\log p(a)
       \left(\log\frac{p(a)}{q(a)}+1\right)\\
 &=\sum_ap(a)\nabla_\theta\log p(a)\log\frac{p(a)}{q(a)},
\end{align*}
最后一步使用 score 均值为零。若精确子集直接求 KL 的梯度，其 $+1$ 项须保留在同样拆分的尾部估计中；若子集直接计算上式的 score 求和，则子集与尾部都按去掉 $+1$ 的形式计算。两种一致拆分都无偏，混合两种形式会遗漏子集概率质量的导数。

集合在求导中固定、尾部校正完整、重要性权重不裁剪，是无偏条件。实际裁剪重要性权重或丢弃尾部会引入偏差。仅说明完整词表 K4 无偏，仍不足以证明它与直接自动微分的精确头部拼接后无偏，因为尾部掩码改变了 score 基线的均值。

\begin{proposition}[精确 KL 头部与 K4 尾部的梯度偏差]
\label{prop:topk-tail-baseline}
固定集合 $A$，精确头部按通常方式自动微分，尾部采用 K4 score 梯度且使用完整未裁剪的重要性校正，则总估计的梯度偏差为
\[
 \E[\widehat{\nabla K}]-\nabla\KL(p\Vert q)
 =\nabla\sum_{a\in A}p(a).
\]
因此，论文 reverse Top-$k$ KL 算法框所写的这种拼接，仅在右侧为零或另有基线校正时才满足梯度无偏。
\end{proposition}
\begin{proof}
记 $s(a)=\nabla\log p(a)$、$l(a)=\log[p(a)/q(a)]$。精确头部梯度为 $\sum_{a\in A}p(a)s(a)[l(a)+1]$，K4 尾部梯度期望为 $\sum_{a\notin A}p(a)s(a)l(a)$。两者相加等于
\[
 \sum_ap(a)s(a)l(a)+\sum_{a\in A}p(a)s(a)
 =\nabla\KL(p\Vert q)+\nabla\sum_{a\in A}p(a).
\]
完整词表的 $\sum_ap(a)s(a)=0$ 不能用在只包含尾部的和式上。
\end{proof}

二分类例子取 $p(1)=\operatorname{sigmoid}(\theta)$、$A=\{1\}$，并在某一点令 $q=p$。真实 KL 梯度为零；K4 尾部也为零，而精确头部梯度为 $p(1)[1-p(1)]>0$。这给出与 teacher 差距为零时仍产生梯度的直接反例。

一种同时保持值与梯度无偏的修正是将尾部 surrogate 写成
\[
 \widehat K_{\mathrm{tail}}^{\mathrm{corr}}
 =\1\{Y\notin A\}\operatorname{stopgrad}\!\left(\frac{p(Y)}{\mu(Y)}\right)
 \left[r_Y\operatorname{stopgrad}(l(Y)+1)-1\right],
 \qquad r_Y=\frac{p(Y)}{\operatorname{stopgrad}(p(Y))}.
\]
在前向值上 $r_Y=1$，括号等于 $l(Y)$，所以值保持原来的尾部估计；自动微分梯度是 $s(Y)[l(Y)+1]$ 乘重要性校正，正好补齐精确头部采用的拆分。另一种修正是头部也用停止梯度的 score surrogate，头尾一起去掉 $+1$。两种选择必须与实现逐项核对。

\subsection{均值递推与均方地板}

\begin{theorem}[固定线性滞后模型的均方界]
设 $D=hI-\alpha\beta F$，$q_D=\|D\|_2<1$。噪声满足
\[
 \E[\xi_t\mid\mathcal F_t]=0,\qquad
 \E[\|\xi_t\|^2\mid\mathcal F_t]\le\sigma^2.
\]
在 $\delta_{t+1}=D\delta_t+\alpha g+\alpha\xi_t$ 中，令 $\delta_*=(I-D)^{-1}\alpha g$，则
\[
 \E\|\delta_t-\delta_*\|^2
 \le q_D^{2t}\E\|\delta_0-\delta_*\|^2
 +\frac{\alpha^2\sigma^2(1-q_D^{2t})}{1-q_D^2}.
\]
\end{theorem}
\begin{proof}
令 $u_t=\delta_t-\delta_*$，稳态方程给出 $u_{t+1}=Du_t+\alpha\xi_t$。条件交叉项为
\[
 \E[\langle Du_t,\xi_t\rangle\mid\mathcal F_t]
 =\langle Du_t,\E[\xi_t\mid\mathcal F_t]\rangle=0.
\]
故 $\E\|u_{t+1}\|^2\le q_D^2\E\|u_t\|^2+\alpha^2\sigma^2$。展开递推、求等比和得到结论。
\end{proof}

一般联合矩阵 $M$ 非正规时，使用 $C_r$、$r>\rho(M)$。输入 $\alpha B\xi_t$、确定初值及鞅差噪声给出
\[
 \E\|z_t\|^2\le C_r^2r^{2t}\|z_0\|^2
 +\frac{\alpha^2C_r^2\|B\|^2\sigma^2(1-r^{2t})}{1-r^2}.
\]
其证明由线性递推展开、不同时间鞅差交叉项为零以及矩阵幂界组成。无偏 Top-$k$ KL 保持固定线性模型的条件均值递推，改变噪声协方差。非线性系统中方差改变未来状态分布，不能扩张成整条随机训练轨迹的无条件均值完全相同。

\section{同步与参数设计的理论结论}

以 teacher tracking 为核心的 MirrorOPD，hard sync 提供独立的 outer reference 和清楚的 exact mirror 收缩率，应作为基准。引入 anchor 平滑、额外 KL 或参数平均后，需要本章的联合系统及事件算子。

$\eta$ 决定基础 mirror target；新增 $\beta$ 将有效步长改为 $\gamma=\eta/(1+\beta\eta)$；$\alpha$ 决定 finite inner 求解速度和离散稳定性；$h$ 决定 anchor 追踪速度；$T_{\mathrm{ema}}$ 与 $\Lambda$ 决定固定 anchor 期间累计的 actor 更新数。这些参数在稳定性上耦合，也不能全部合并为一个有效学习率。

检查应与实现时序一致：逐 actor step 更新前 EMA 使用定理\ref{thm:ema-jury}或推论\ref{cor:ema-opd-jury}；多个 inner steps 后同步使用 $SK^N$；logit 几何同步使用定理\ref{thm:ema-geometric-exact}。参数 EMA 还需检查局部近似误差、状态测度变化与 teacher 不可表示的影响。

EMA 的 exact anchor 推进比 hard sync 慢，已由精确收缩因子证明。实际价值应通过稳定预算与噪声预算评估：是否允许更大的 $\alpha$，是否减小训练波动，是否缓解 reference 突变。改善需要额外条件或实验证据；lag 稳定定理本身只说明 actor 与 anchor 的距离如何演化。
